\chapter{Evaluación}
\label{sec:paper-eval}

La evaluación se organiza en tres ejes: (i) benchmarks
cuantitativos de latencia y throughput, (ii) property-based tests
que verifican invariantes, (iii) smoke end-to-end sobre datasets
sintéticos. Cada eje responde a una pregunta distinta.

\section{Benchmarks de performance}

Los benchmarks viven en
\filepath{core/graph-engine/benches/} y se corren con
\emph{criterion}. Reportamos números medidos con \emph{baseline
\code{pre-l}} (cierre del \emph{M-track} de abril 2026) sobre
\emph{Galaxy Book4} (Intel Core Ultra hybrid P/E, \emph{thin LTO},
sin \code{target-cpu=native} --- ese flag regresiona en CPUs
híbridas, ver \filepath{docs/perf/release-profile.md}). Mediana
sobre 100 muestras criterion, $2$~s \emph{warm-up}, $5$~s
\emph{measurement}.

\begin{table}[H]
\centering
\small
\begin{tabular}{@{}lrrr@{}}
\toprule
Benchmark & \emph{Pre-LTO} & \emph{Pre-L} & \emph{Mejora} \\
\midrule
Hunt 1-hop, $|V|=10^3$   & 302 µs   & 97.6 µs  & $3.1\times$ \\
Hunt 1-hop, $|V|=10^4$   & 3.4 ms   & 874 µs   & $3.9\times$ \\
Hunt 1-hop, $|V|=10^5$   & 40.2 ms  & \textbf{13.5 ms} & $3.0\times$ \\
Match 2-hop, $|V|=10^4$  & ---      & 4.46 ms  & --- \\
Match 3-hop, $|V|=10^4$  & ---      & 10.23 ms & --- \\
Match 4-hop, $|V|=10^4$  & 56 ms    & \textbf{20.36 ms} & $2.7\times$ \\
Ingest sysmon NDJSON, 1k & 5.4 ms   & 4.74 ms (210K ev/s) & $1.1\times$ \\
Ingest sysmon NDJSON, 10k & 66 ms   & 38.79 ms (250K ev/s) & $1.7\times$ \\
Ingest sysmon NDJSON, 50k & 327 ms  & 204.4 ms (240K ev/s) & $1.6\times$ \\
Score path (depth 4) & 1.04 µs    & 401 ns   & $2.6\times$ \\
\bottomrule
\end{tabular}
\caption{Latencia y \emph{throughput} sobre las cinco \emph{suites}
canónicas \emph{matcher} y la \emph{suite} de \emph{ingest}. La
columna \emph{pre-LTO} es el \emph{baseline} previo a la
optimización de \emph{profile} y al \emph{M-track} (\emph{S3, M1,
M2, M3} y refinamientos M3$+$/M3$++$); la columna \emph{pre-L} es
el estado al cierre de abril 2026. Datos completos
en \filepath{docs/paper/data/bench-table.md}, generado por
\filepath{scripts/bench-report.sh}.}
\label{tab:paper-bench}
\end{table}

\textbf{Operadores canónicos WCOJ.} El \emph{M-track} agregó cuatro
\emph{count operators} con \emph{symmetry-breaking} \emph{hardcoded}
para las \emph{shapes} estructuralmente conocidas (\(K_3\), \(K_4\),
\(K_5\), \(C_6\)). En el extremo denso del \emph{sweep} \emph{ER}
(\(V=1000\), \(E=9913\), $C_6=926763$ \emph{instances}), \emph{naive}
\(O(m \cdot d^4 \cdot \log d)\) tarda \(2.03\)~s, mientras que la
\emph{canonical+par} con \emph{star-from-min} \(a < \{b,c,d,e,f\}\)
sobre \(\mathbb{Z}/6\) tarda $22.5$~ms --- \(\sim 90\times\). Para
\(K_3/K_4/K_5\) el corte respectivo es \(11\text{--}12\times\) sobre
\emph{naive} con \emph{chain} \(a < b < c[<d[<e]]\). Detalle en
\filepath{docs/perf/latex/cambios-perf-2026-04.tex} (capítulos
M3 y M3$++$).

Observaciones:

\begin{itemize}
    \item Hunt 1-hop a $10^5$ nodos en \(13.5\)~ms ronda el
    \emph{abstract claim} \emph{sub-10 ms} de \emph{Black Hat
    Arsenal 2026} con un \emph{gap} restante de \(\sim 35\%\)
    (cubrible por la C-tier \emph{frontier-batched + AMAC}, ya
    presente bajo \emph{feature flag}).
    \item \emph{Ingest} a $250$K\,ev/s sigue $4\times$ bajo el
    \emph{target} de $1$M\,ev/s. La \emph{lift} restante depende
    de \emph{simd-json} con \emph{schema-known fast-path},
    no incluido en el \emph{M-track}.
    \item La composición de score es barata: escala linealmente
    con la longitud del path. La \emph{depth 4} cae $2.6\times$
    por la migración \code{HuntResult: Vec<String>} \(\to\)
    \code{Vec<Arc<str>>}, que evita
    \(\text{path.len()} \cdot \text{String::from(\dots)}\)
    \emph{allocations} por \emph{result}.
\end{itemize}

\section{Property-based tests}

Tres propiedades clave se testan con
\emph{proptest}:

\begin{enumerate}
    \item \textbf{Equivalencia DFS$\leftrightarrow$SIMD}.
    \filepath{core/matcher-ffi/tests/equivalence_dfs_simd.rs}
    genera grafos aleatorios con distribución de grados power-law
    y patrones de longitud 2-4, y verifica
    $\text{SIMD}(G, h) = \text{DFS}(G, h)$ como conjuntos. $10^4$
    casos por corrida en CI; $10^5$ en pass local; ningún fallo
    reportado.
    \item \textbf{DSL roundtrip}.
    \filepath{platform/dsl/tests/roundtrip.rs} genera ASTs
    aleatorios, los imprime como texto DSL, los re-parsea, y
    verifica igualdad. $10^4$ casos por corrida.
    \item \textbf{Estabilidad numérica del score}.
    \filepath{core/graph-engine/tests/anomaly_proptest.rs} perturba
    el grafo (agrega/remueve una arista) y verifica que los
    scores en el nuevo grafo difieren por menos de una cota
    proporcional al tamaño de la perturbación.
\end{enumerate}

Estos tests capturan clases enteras de bugs que no caen en el
radar de tests unitarios con ejemplos fijos.

\section{Smoke end-to-end manual}

Sobre los datasets de
\filepath{demo_data/} (ver libro \emph{Manual}, apéndice C) se
verifica el flujo completo en cada release candidate:

\begin{enumerate}
    \item Crear sesión, ingerir
    \filepath{fortigate_sample.csv}. Verificar: 50k eventos
    procesados, $\le$50 en DLQ, schema visible en catálogo.
    \item Correr
    \code{IP -[Connect]-> IP -[Connect]-> IP {k=3, window=30m}}.
    Verificar: $\ge$3 paths con score $>$0.6,
    \code{explain\_hunt\_result} devuelve los 5 componentes.
    \item Expandir un nodo sospechoso. Verificar pestañas timeline,
    events, properties pobladas.
    \item Taggear path como IOC. Exportar IOCs en formato Sentinel.
    Verificar JSON con estructura \code{ThreatIntelligenceIndicator}.
    \item Cerrar y reabrir la sesión. Verificar: todo el estado
    persistido (grafo, tags, notas, hunts históricos).
\end{enumerate}

El smoke está documentado en
\filepath{docs/_archive/phase1/MIGRATION_PLAN.md}; sobrevive
como checklist operativo de release.

\section{Métricas del proyecto}

Tamaño y complejidad del codebase:

\begin{table}[H]
\centering
\small
\begin{tabular}{@{}lr@{}}
\toprule
Métrica & Valor \\
\midrule
Crates Rust en \code{core/}      & 4 \\
Crates Rust en \code{platform/}  & 9 \\
Crates Rust en \code{apps/}      & 3 \\
Crate TypeScript (MCP)           & 1 (55 tools, 10 dominios) \\
Biblioteca C++ (libgraphmatch)   & 1 (ABI v1) \\
Tests unitarios + integración     & $\sim$450 \\
Property-based tests             & 3 suites (equiv, DSL, anomaly) \\
Benchmarks                       & 4 (hunt, dedup, score, ingest) \\
ADRs                             & 7 \\
\bottomrule
\end{tabular}
\caption{Inventario del proyecto a fecha de corte.}
\label{tab:paper-metricas}
\end{table}

\section{Comparación pre/post shearing-layers}

Pre-rediseño, el grafo de dependencias tenía 19 aristas de
``capas'' cruzando sentidos mezclados, con al menos dos ciclos
frágiles entre parsers y el motor de matching. Post-rediseño hay
0 ciclos y 2 aristas residuales \code{core}$\to$\code{platform}
(ver cap.~\ref{sec:paper-limites}). El
\emph{cargo check} del workspace pasó de $\sim$2m15s a $\sim$1m05s
en frío en el mismo hardware.

El costo directo del rediseño fue un sprint de tres semanas
(Fases 2+3+4 cerradas). El costo evitado es difícil de
cuantificar: cada feature nueva que no toca el core se paga sola
en ahorro de builds y revisiones.
